\documentclass{beamer}

\usetheme{Berlin}
\usecolortheme{Orchid}
\useoutertheme{miniframes}
\setbeamertemplate{navigation symbols}{}

\usepackage[utf8]{inputenc}
\usepackage[T1]{fontenc}
\usepackage{amsmath}
\usepackage{amssymb}
\usepackage{booktabs}
\usepackage{graphicx}
\graphicspath{{../images/}}

\usepackage{xcolor}
\usepackage{listings}
\usepackage{hyperref}
\usepackage{tikz}
\usetikzlibrary{arrows.meta,positioning,shapes.geometric,calc}

\lstset{
  language=Java,
  basicstyle=\ttfamily\small,
  keywordstyle=\color{blue},
  commentstyle=\color{gray},
  stringstyle=\color{teal},
  showstringspaces=false,
  columns=fullflexible,
  breaklines=true
}

% Per-slide source line (references shown on the slide itself).
\newcommand{\slidesources}[1]{%
  \vfill
  {\tiny\color{gray}\raggedright\sloppy\parbox{\linewidth}{Sources: #1}\par}%
}

% Unit-wise source strings for this subject (used by slides/notes).
% Path is relative to the lecture `latex/` folder (the compilation CWD).
% OOPS with Java (CSEG1044) — unit-wise sources
%
% These are short, student-facing reference strings intended to be used with:
%   - \slidesources{...}  (in beamer slides)
%   - \notesources{...}   (in lecture notes)
%
% Style inspiration: other subjects in My-Subjects/ use semicolon-separated sources
% and include an "accessed" date for web documentation.

\newcommand{\OOPSUnitISources}{Schildt, \textit{Java: A Beginner's Guide} (10e), McGraw-Hill (Java fundamentals + OOP basics).; Oracle Java Tutorials: Getting Started + Language Basics + Classes and Objects (accessed \today).; Java SE API docs (e.g., \texttt{String}, \texttt{Scanner}) (accessed \today).}

\newcommand{\OOPSUnitIISources}{Schildt, \textit{Java: The Complete Reference} (12e), McGraw-Hill (classes, inheritance, packages, interfaces).; Bloch, \textit{Effective Java} (3e), Addison-Wesley, 2018 (design choices; interfaces vs abstract classes).; Oracle Java Tutorials: Classes and Objects + Interfaces + Packages (accessed \today).; Java SE API docs (e.g., \texttt{Comparable}, \texttt{Runnable}) (accessed \today).}

\newcommand{\OOPSUnitIIISources}{Schildt, \textit{Java: The Complete Reference} (12e) (nested/inner classes, exceptions, threads, I/O).; Oracle Java Tutorials: Nested Classes + Exceptions + Concurrency + Basic I/O (accessed \today).; Java SE API docs (e.g., \texttt{Thread}, \texttt{AutoCloseable}) (accessed \today).}

\newcommand{\OOPSUnitIVSources}{Schildt, \textit{Java: The Complete Reference} (12e) (generics, lambdas, Swing, JDBC).; Bloch, \textit{Effective Java} (3e) (generics + lambdas).; Oracle Java Tutorials: Generics + Lambda Expressions + Swing + JDBC Basics (accessed \today).; Java SE API docs (e.g., \texttt{java.util.function}, \texttt{javax.swing}, \texttt{java.sql}) (accessed \today).}

\newcommand{\OOPSUnitVSources}{Schildt, \textit{Java: The Complete Reference} (12e) (Collections Framework + wrapper classes).; Oracle Java docs: Collections Framework overview (accessed \today).; Java SE API docs (\texttt{java.util} collections; wrapper classes in \texttt{java.lang}) (accessed \today).}

\newcommand{\OOPSUnitVISources}{Git SCM: \textit{Pro Git} (2e) (version control workflow), accessed \today.; GitHub Docs (repositories, issues, pull requests), accessed \today.; Oracle Java Tutorials: Swing + JDBC (accessed \today).; JUnit 5 User Guide (accessed \today).; Apache Maven Project: Getting Started (accessed \today).; Gradle User Manual: Getting Started (accessed \today).; UML class diagrams: UML 2.x overview/spec (accessed \today).}

\newcommand{\OOPSMiscWhyInterfacesSources}{Schildt, \textit{Java: The Complete Reference} (12e) (interfaces).; Bloch, \textit{Effective Java} (3e) (prefer interfaces where appropriate).; Oracle Java Tutorials: Interfaces (accessed \today).; Java SE API docs (e.g., \texttt{Comparable}, \texttt{Runnable}) (accessed \today).}



\title[OOPS with Java]{Object Oriented Programming (CSEG1044)}
\subtitle{Unit 02 -- Lecture 05: Abstract Classes, `final`, and Interfaces (Design Choices)}
\begin{document}

\begin{frame}[fragile]
  \titlepage
  \vspace{0.5em}
  \slidesources{\OOPSUnitIISources}
\end{frame}

\begin{frame}{Quick Links}
  \centering
  \hyperlink{sec:overview}{\beamerbutton{Overview}}\hspace{0.6em}
  \hyperlink{sec:concepts}{\beamerbutton{Key Topics}}\hspace{0.6em}
  \hyperlink{sec:demo}{\beamerbutton{Demo}}\hspace{0.6em}
  \hyperlink{sec:summary}{\beamerbutton{Summary}}
  \slidesources{\OOPSUnitIISources}
\end{frame}

\begin{frame}{Agenda}
  \tableofcontents
  \slidesources{\OOPSUnitIISources}
\end{frame}

\section{Overview}
\label{sec:overview}

\begin{frame}{Lecture Snapshot}
\begin{itemize}[<+->]
  \item \textbf{Unit focus:} Classes, Inheritance, Packages, and Interfaces
  \item \textbf{Course thread:} Use Abstract Classes, `final`, and Interfaces (Design Choices) to strengthen the domain model, APIs, and access rules inside Campus Resource Manager.
\end{itemize}
  \slidesources{\OOPSUnitIISources}
\end{frame}

\begin{frame}{Recall Before You Start}
\begin{itemize}[<+->]
  \item \textbf{Inheritance}: Abstract classes extend inheritance by allowing an incomplete but shared parent design. Main reference: \texttt{Unit 02 / Lecture 04 slides/notes}.
  \item \textbf{Interface idea}: You already saw interfaces as contracts; this lecture compares them with abstract classes in more depth. Main reference: \texttt{Unit 02 / Lecture 03 slides/notes}.
\end{itemize}
  \slidesources{\OOPSUnitIISources}
\end{frame}


\begin{frame}{Learning Outcomes}
  \begin{itemize}[<+->]
    \item Use abstract classes for shared code + required methods.
    \item Explain \texttt{final} for variables/methods/classes.
    \item Use interfaces (including \texttt{Comparable}, \texttt{Runnable}) as contracts.
  \end{itemize}
  \slidesources{\OOPSUnitIISources}
\end{frame}

\begin{frame}{Key Terms to Know}
\begin{itemize}[<+->]
  \item \textbf{Abstract class}: class that cannot be instantiated; may include abstract methods.
  \item \textbf{Abstract method}: declared without body; must be implemented by concrete subclasses.
  \item \textbf{Interface}: contract specifying methods a class must provide.
  \item \textbf{Final}: ``cannot change'' (variable), ``cannot override'' (method), ``cannot extend'' (class).
  \item \textbf{Contract}: what is promised by a type (what callers can rely on).
\end{itemize}
  \slidesources{\OOPSUnitIISources}
\end{frame}

\begin{frame}{Study Flow for This Lecture}
\begin{enumerate}[<+->]
  \item Refresh the prerequisite bridge before reading new ideas in isolation.
  \item Study the main build in this order: \textit{Abstract classes and methods} $\rightarrow$ \textit{\texttt{final} keyword (design intent)} $\rightarrow$ \textit{Interfaces: built-in + user-defined}.
  \item Trace the worked example and connect each line back to one concept frame.
  \item Attempt the mini exercises before moving to the recap and exit check.
\end{enumerate}
  \slidesources{\OOPSUnitIISources}
\end{frame}


\section{Key Topics}
\label{sec:concepts}

\begin{frame}{Unit 02: Classes, Inheritance, Packages and Interfaces}
  \begin{itemize}[<+->]
    \item Lecture focus: Abstract Classes, `final`, and Interfaces (Design Choices)
  \end{itemize}
  \slidesources{\OOPSUnitIISources}
\end{frame}

\begin{frame}{Today: Roadmap}
  \begin{itemize}[<+->]
    \item Abstract classes/methods; partial implementations
    \item `final` variables/methods/classes (design intent)
    \item Interfaces: built-in examples (`Comparable`, `Runnable`) + user-defined interfaces
  \end{itemize}
  \slidesources{\OOPSUnitIISources}
\end{frame}

\begin{frame}{Abstract class (when to use)}
  \begin{itemize}[<+->]
    \item Use when you want to share code + force subclasses to implement some methods.
    \item Abstract class can have fields + concrete methods + abstract methods.
    \item Cannot instantiate: \texttt{new Shape()} if \texttt{Shape} is abstract.
  \end{itemize}
  \slidesources{\OOPSUnitIISources}
\end{frame}

\begin{frame}{\texttt{final} keyword (meaning)}
  \begin{itemize}[<+->]
    \item \textbf{final variable}: assigned once (constant-like).
    \item \textbf{final method}: cannot be overridden.
    \item \textbf{final class}: cannot be extended (e.g., \texttt{String}).
  \end{itemize}
  \slidesources{\OOPSUnitIISources}
\end{frame}

\begin{frame}{Interface (when to use)}
  \begin{itemize}[<+->]
    \item Use when you want a contract with multiple implementations.
    \item Reduces coupling: code depends on interface, not concrete class.
    \item A class can implement multiple interfaces.
  \end{itemize}
  \slidesources{\OOPSUnitIISources}
\end{frame}

\begin{frame}{Built-in interfaces (examples)}
  \begin{itemize}[<+->]
    \item \texttt{Comparable<T>}: defines natural ordering (\texttt{compareTo}).
    \item \texttt{Runnable}: defines a task (\texttt{run}) to execute in a thread.
  \end{itemize}
  \slidesources{\OOPSUnitIISources}
\end{frame}

\begin{frame}{Checkpoint and Reflection}
\begin{block}{Reflection prompt}
Where would \textit{Abstract Classes, `final`, and Interfaces (Design Choices)} matter most in Campus Resource Manager, and which design choice would you justify using this lecture's ideas?
\end{block}
\begin{block}{Quick self-check}
Which part needs one more example before you move on: \textit{Abstract classes and methods}, \textit{\texttt{final} keyword (design intent)}, the worked example, or the recap?
\end{block}
  \slidesources{\OOPSUnitIISources}
\end{frame}

\section{Demo / Example}
\label{sec:demo}

\begin{frame}[fragile]{Example: \texttt{Comparable} for sorting}
\begin{lstlisting}
import java.util.Arrays;

class Student implements Comparable<Student> {
    private final String rollNo;
    private final double cgpa;
    Student(String rollNo, double cgpa) { this.rollNo = rollNo; this.cgpa = cgpa; }
    @Override public int compareTo(Student o) {
        return Double.compare(o.cgpa, this.cgpa); // descending
    }
}
\end{lstlisting}
  \slidesources{\OOPSUnitIISources}
\end{frame}

\begin{frame}[fragile]{Demo usage}
\begin{lstlisting}
public class Demo {
    public static void main(String[] args) {
        Student[] a = {
            new Student("21CS1001", 8.2),
            new Student("21CS1002", 9.1)
        };
        Arrays.sort(a); // uses compareTo
        System.out.println(Arrays.toString(a));
    }
}
\end{lstlisting}
  \slidesources{\OOPSUnitIISources}
\end{frame}

\begin{frame}{Mini Exercises (2--3 mins)}
  \begin{enumerate}[<+->]
    \item Can we instantiate an abstract class?
    \item What does \texttt{final class} prevent?
    \item One sentence: why use interfaces?
  \end{enumerate}
  \slidesources{\OOPSUnitIISources}
\end{frame}

\begin{frame}{Watch-outs}
\begin{itemize}[<+->]
  \item Trying to instantiate an abstract class directly.
  \item Marking too many classes or methods final too early.
  \item Choosing an interface even when shared state or shared code is central.
\end{itemize}
  \slidesources{\OOPSUnitIISources}
\end{frame}

\section{Summary}
\label{sec:summary}

\begin{frame}{Key Takeaways}
  \begin{itemize}[<+->]
    \item Abstract class: share code + require methods.
    \item Interface: contract + low coupling + multiple implementations.
    \item \texttt{final} communicates ``no change'' / ``no override'' / ``no extend''.
  \end{itemize}
  \slidesources{\OOPSUnitIISources}
\end{frame}

\begin{frame}{Checkpoint Questions}
  \begin{enumerate}[<+->]
    \item When do you prefer interface over abstract class?
    \item Give one use case for \texttt{final method}.
    \item What does \texttt{Comparable} enable?
  \end{enumerate}
  \slidesources{\OOPSUnitIISources}
\end{frame}

\begin{frame}{Next Study Steps}
\begin{itemize}[<+->]
  \item Revisit the recall references if any older idea still feels weak.
  \item Rework the lecture example without looking at the notes first.
  \item Attempt the self-study practice and then answer the exit questions in full sentences.
  \item Apply the idea once inside Campus Resource Manager to make the concept stick.
\end{itemize}
  \slidesources{\OOPSUnitIISources}
\end{frame}

\begin{frame}{Next Unit Preview}
  \textbf{Next:} Unit 03 -- Nested Classes, Exceptions, Multithreading, I/O Streams
  \vspace{0.6em}
  \begin{itemize}[<+->]
    \item We will start with nested classes and exception handling.
  \end{itemize}
  \slidesources{\OOPSUnitIISources}
\end{frame}

\end{document}
